% ============================================================
\chapter{Pipelines de Donn\'ees et Feature Stores}
\label{ch:pipelines}
\index{pipelines de donn\'ees}
% ============================================================

\section{Qu'est-ce qu'un pipeline de donn\'ees~?}
\label{sec:data-pipeline-def}

Un \textbf{pipeline de donn\'ees}\index{pipeline de donn\'ees} est une
s\'equence automatis\'ee d'\'etapes qui transforme des donn\'ees brutes en
donn\'ees exploitables par un mod\`ele de machine learning. Dans un contexte
MLOps, les pipelines assurent la \emph{reproductibilit\'e}, la
\emph{tra\c{c}abilit\'e} et l'\emph{automatisation} du flux de donn\'ees.

\begin{definition}[Pipeline de donn\'ees]
Un pipeline de donn\'ees est un ensemble d'\'etapes de traitement
orchestr\'ees, o\`u chaque \'etape consomme des entr\'ees, effectue une
transformation, et produit des sorties. Les \'etapes sont
\textbf{id\'empotentes} (ex\'ecuter deux fois donne le m\^eme r\'esultat)
et \textbf{d\'eterministes} (m\^emes entr\'ees $\Rightarrow$ m\^emes
sorties).
\end{definition}

\subsection{ETL vs ELT}
\index{ETL}\index{ELT}

Deux paradigmes classiques structurent les pipelines de donn\'ees~:

\begin{center}
\begin{tabular}{@{}lp{5.5cm}p{5.5cm}@{}}
\toprule
& \textbf{ETL} & \textbf{ELT} \\
\midrule
\textbf{Principe} & Extract $\to$ Transform $\to$ Load &
  Extract $\to$ Load $\to$ Transform \\
\textbf{Transformation} & Avant le chargement (en m\'emoire) &
  Apr\`es le chargement (dans le data warehouse) \\
\textbf{Outils} & Airflow, Prefect, Luigi &
  dbt, Spark SQL, BigQuery \\
\textbf{Cas d'usage} & Donn\'ees structur\'ees, volume mod\'er\'e &
  Grands volumes, stockage bon march\'e \\
\bottomrule
\end{tabular}
\end{center}

\begin{remarque}
En MLOps, on utilise souvent un hybride~: ELT pour la pr\'eparation
g\'en\'erale des donn\'ees, puis un pipeline ML sp\'ecifique (feature
engineering, entra\^inement, \'evaluation) orchestr\'e s\'epar\'ement.
\end{remarque}

\section{Outils d'orchestration}
\label{sec:orchestration-tools}
\index{orchestration}

\subsection{Apache Airflow}
\index{Airflow}

\textbf{Airflow} est l'orchestrateur le plus r\'epandu. Il utilise des
\textbf{DAGs} (graphes orient\'es acycliques) d\'efinis en Python pour
d\'ecrire les d\'ependances entre t\^aches.

\begin{codeblock}[title=\textbf{Exemple de DAG Airflow}]
\begin{minted}{python}
from airflow import DAG
from airflow.operators.python import PythonOperator
from datetime import datetime

def extract():
    """Extract data from source."""
    print("Extracting data...")

def transform():
    """Transform raw data."""
    print("Transforming data...")

def load():
    """Load data to destination."""
    print("Loading data...")

with DAG(
    dag_id="ml_data_pipeline",
    start_date=datetime(2025, 1, 1),
    schedule_interval="@daily",
    catchup=False,
) as dag:
    t1 = PythonOperator(task_id="extract", python_callable=extract)
    t2 = PythonOperator(task_id="transform", python_callable=transform)
    t3 = PythonOperator(task_id="load", python_callable=load)

    t1 >> t2 >> t3  # dependances
\end{minted}
\end{codeblock}

\subsection{Prefect}
\index{Prefect}

\textbf{Prefect} adopte une approche plus pythonique avec des d\'ecorateurs
\cli{@flow} et \cli{@task}, sans la complexit\'e de configuration d'Airflow.

\subsection{Luigi}
\index{Luigi}

\textbf{Luigi} (par Spotify) est orient\'e cibles~: chaque t\^ache d\'efinit
une sortie (\cli{output()}) et ses d\'ependances (\cli{requires()}). Le
graphe est construit automatiquement.

\section{\'Ecrire un pipeline avec Prefect}
\label{sec:prefect-pipeline}

\begin{codeblock}[title=\textbf{Pipeline ML complet avec Prefect}]
\begin{minted}{python}
from prefect import flow, task
import pandas as pd
from sklearn.model_selection import train_test_split
from sklearn.ensemble import RandomForestClassifier
from sklearn.metrics import accuracy_score, f1_score
import joblib
from pathlib import Path


@task(retries=2, retry_delay_seconds=10)
def extract_data(path: str) -> pd.DataFrame:
    """Load raw data from CSV."""
    df = pd.read_csv(path)
    print(f"Loaded {len(df)} rows from {path}")
    return df


@task
def clean_data(df: pd.DataFrame) -> pd.DataFrame:
    """Remove duplicates and handle missing values."""
    df = df.drop_duplicates()
    df = df.dropna(subset=["target"])
    df = df.fillna(df.median(numeric_only=True))
    print(f"After cleaning: {len(df)} rows")
    return df


@task
def feature_engineering(df: pd.DataFrame) -> pd.DataFrame:
    """Create new features."""
    numeric_cols = df.select_dtypes(include="number").columns
    for col in numeric_cols:
        if col != "target":
            df[f"{col}_squared"] = df[col] ** 2
            df[f"{col}_log"] = df[col].clip(lower=1e-6).apply(
                lambda x: x ** 0.5
            )
    return df


@task
def split(df: pd.DataFrame, test_size: float = 0.2):
    """Split data into train and test sets."""
    X = df.drop("target", axis=1)
    y = df["target"]
    return train_test_split(X, y, test_size=test_size, random_state=42)


@task
def train_model(X_train, y_train, n_estimators: int = 100):
    """Train a RandomForest classifier."""
    model = RandomForestClassifier(
        n_estimators=n_estimators, random_state=42, n_jobs=-1
    )
    model.fit(X_train, y_train)
    return model


@task
def evaluate_model(model, X_test, y_test) -> dict:
    """Evaluate model and return metrics."""
    y_pred = model.predict(X_test)
    metrics = {
        "accuracy": accuracy_score(y_test, y_pred),
        "f1_score": f1_score(y_test, y_pred, average="weighted"),
    }
    print(f"Metrics: {metrics}")
    return metrics


@task
def save_model(model, path: str) -> None:
    """Save trained model to disk."""
    Path(path).parent.mkdir(parents=True, exist_ok=True)
    joblib.dump(model, path)
    print(f"Model saved to {path}")


@flow(name="ml-training-pipeline")
def training_pipeline(
    data_path: str = "data/dataset.csv",
    model_path: str = "models/model.joblib",
    n_estimators: int = 100,
):
    """End-to-end ML training pipeline."""
    # Extract
    raw_df = extract_data(data_path)

    # Transform
    clean_df = clean_data(raw_df)
    feat_df = feature_engineering(clean_df)

    # Split and train
    X_train, X_test, y_train, y_test = split(feat_df)
    model = train_model(X_train, y_train, n_estimators)

    # Evaluate and save
    metrics = evaluate_model(model, X_test, y_test)
    save_model(model, model_path)

    return metrics


if __name__ == "__main__":
    training_pipeline()
\end{minted}
\end{codeblock}

\begin{codeblock}[title=\textbf{Ex\'ecution et d\'eploiement Prefect}]
\begin{minted}{bash}
# Installer Prefect
pip install prefect

# Executer localement
python pipeline.py

# Demarrer le serveur Prefect (UI)
prefect server start

# Deployer le flow
prefect deploy pipeline.py:training_pipeline \
    --name "ml-training" \
    --cron "0 6 * * *"  # tous les jours a 6h
\end{minted}
\end{codeblock}

\section{DVC Pipelines~: orchestrer un pipeline ML}
\label{sec:dvc-pipelines}
\index{DVC!pipelines}

\textbf{DVC}\index{DVC} (Data Version Control) permet de d\'efinir des
pipelines ML dans un fichier \file{dvc.yaml}. Chaque \'etape sp\'ecifie
ses d\'ependances et ses sorties, et DVC ne r\'e-ex\'ecute que les \'etapes
dont les entr\'ees ont chang\'e.

\begin{codeblock}[title=\textbf{dvc.yaml --- Pipeline ML}]
\begin{minted}{yaml}
stages:
  prepare:
    cmd: python src/prepare.py --input data/raw.csv
         --output data/prepared.csv
    deps:
      - src/prepare.py
      - data/raw.csv
    outs:
      - data/prepared.csv

  featurize:
    cmd: python src/featurize.py --input data/prepared.csv
         --output data/features.csv
    deps:
      - src/featurize.py
      - data/prepared.csv
    outs:
      - data/features.csv

  train:
    cmd: python src/train.py --input data/features.csv
         --model models/model.pkl
    deps:
      - src/train.py
      - data/features.csv
    outs:
      - models/model.pkl
    params:
      - train.n_estimators
      - train.max_depth

  evaluate:
    cmd: python src/evaluate.py --model models/model.pkl
         --data data/features.csv
    deps:
      - src/evaluate.py
      - models/model.pkl
      - data/features.csv
    metrics:
      - metrics.json:
          cache: false
    plots:
      - plots/confusion_matrix.png
\end{minted}
\end{codeblock}

\begin{codeblock}[title=\textbf{Commandes DVC pipelines}]
\begin{minted}{bash}
# Executer le pipeline complet
dvc repro

# Visualiser le graphe du pipeline
dvc dag

# Voir les metriques
dvc metrics show

# Comparer avec une autre branche
dvc metrics diff main
\end{minted}
\end{codeblock}

\begin{bonnepratique}[title=\textbf{Quand utiliser DVC vs Prefect/Airflow}]
\begin{itemize}
  \item \textbf{DVC}~: pipelines ML li\'es au code (m\^eme d\'ep\^ot Git),
    ex\'ecution locale ou CI/CD, versionnement des donn\'ees int\'egr\'e.
  \item \textbf{Prefect/Airflow}~: orchestration de flux complexes,
    ex\'ecution planifi\'ee, d\'ependances entre syst\`emes (API, bases de
    donn\'ees), monitoring avanc\'e.
\end{itemize}
\end{bonnepratique}

\section{Feature Stores}
\label{sec:feature-stores}
\index{feature store}

\subsection{Pourquoi un feature store~?}

Le \textbf{feature engineering}\index{feature engineering} est souvent la
partie la plus co\^uteuse d'un projet ML. Sans feature store~:
\begin{itemize}
  \item Les features sont recalcul\'ees \`a chaque exp\'erience
  \item Le code de feature engineering est dupliqu\'e entre \'equipes
  \item Les features d'entra\^inement et de serving diff\`erent
    (\emph{training-serving skew}\index{training-serving skew})
  \item Pas de tra\c{c}abilit\'e sur l'origine des features
\end{itemize}

\begin{definition}[Feature Store]
Un \textbf{feature store} est une plateforme centralis\'ee pour stocker,
g\'erer, d\'ecouvrir et servir des features ML. Il garantit la
\textbf{coh\'erence} entre l'entra\^inement (acc\`es \emph{batch}) et
l'inf\'erence (acc\`es \emph{temps r\'eel}).
\end{definition}

\subsection{Architecture d'un feature store}

\begin{center}
\begin{tikzpicture}[
  box/.style={rectangle, draw, rounded corners, minimum width=2.8cm,
    minimum height=1cm, align=center, font=\small},
  store/.style={rectangle, draw, rounded corners, minimum width=3.5cm,
    minimum height=2.5cm, align=center, font=\small, fill=blue!8},
  arr/.style={-{Stealth[length=2.5mm]}, thick},
  data/.style={cylinder, draw, shape border rotate=90, aspect=0.25,
    minimum width=1.5cm, minimum height=1cm, font=\scriptsize, fill=gray!10}
]
  % Sources
  \node[box, fill=green!10] (src1) {Base de\\donn\'ees};
  \node[box, fill=green!10, below=0.6cm of src1] (src2) {Fichiers\\CSV/Parquet};
  \node[box, fill=green!10, below=0.6cm of src2] (src3) {Flux\\temps r\'eel};

  % Feature Store
  \node[store, right=2.5cm of src2] (fs) {};
  \node[above=0.05cm, font=\small\bfseries] at (fs.north) {Feature Store};
  \node[data] at ([yshift=0.2cm]fs.center) {Offline};
  \node[data] at ([yshift=-0.8cm]fs.center) {Online};

  % Consumers
  \node[box, fill=orange!10, right=2.5cm of fs, yshift=0.6cm] (train)
    {Entra\^inement};
  \node[box, fill=orange!10, right=2.5cm of fs, yshift=-0.6cm] (serve)
    {Serving};

  % Arrows
  \draw[arr] (src1.east) -- (fs.west|-src1);
  \draw[arr] (src2.east) -- (fs.west);
  \draw[arr] (src3.east) -- (fs.west|-src3);
  \draw[arr] (fs.east|-train) -- (train.west);
  \draw[arr] (fs.east|-serve) -- (serve.west);

  % Labels
  \node[font=\scriptsize, above=0.1cm of src1, text=gray] {Sources};
  \node[font=\scriptsize, above=0.1cm of train, text=gray] {Consommateurs};
\end{tikzpicture}
\end{center}

\subsection{Feast~: un feature store open-source}
\index{Feast}

\textbf{Feast}\index{Feast} est le feature store open-source le plus
populaire. Il g\`ere les d\'efinitions de features, le stockage offline
(fichiers Parquet, BigQuery) et online (Redis, DynamoDB).

\begin{codeblock}[title=\textbf{D\'efinition de features avec Feast}]
\begin{minted}{python}
# feature_repo/features.py
from feast import Entity, FeatureView, Field, FileSource
from feast.types import Float32, Int64
from datetime import timedelta

# Source de donnees
customer_source = FileSource(
    path="data/customer_features.parquet",
    timestamp_field="event_timestamp",
)

# Entite
customer = Entity(
    name="customer_id",
    join_keys=["customer_id"],
)

# Vue de features
customer_features = FeatureView(
    name="customer_features",
    entities=[customer],
    ttl=timedelta(days=1),
    schema=[
        Field(name="total_purchases", dtype=Int64),
        Field(name="avg_order_value", dtype=Float32),
        Field(name="days_since_last_order", dtype=Int64),
        Field(name="purchase_frequency", dtype=Float32),
    ],
    source=customer_source,
    online=True,
)
\end{minted}
\end{codeblock}

\begin{codeblock}[title=\textbf{Utilisation de Feast}]
\begin{minted}{python}
from feast import FeatureStore
import pandas as pd

store = FeatureStore(repo_path="feature_repo/")

# Recuperer des features pour l'entrainement (offline)
entity_df = pd.DataFrame({
    "customer_id": [1001, 1002, 1003],
    "event_timestamp": pd.to_datetime(["2025-01-15"] * 3),
})
training_df = store.get_historical_features(
    entity_df=entity_df,
    features=[
        "customer_features:total_purchases",
        "customer_features:avg_order_value",
        "customer_features:purchase_frequency",
    ],
).to_df()

# Recuperer des features en temps reel (online)
online_features = store.get_online_features(
    features=[
        "customer_features:total_purchases",
        "customer_features:avg_order_value",
    ],
    entity_rows=[{"customer_id": 1001}],
).to_dict()
\end{minted}
\end{codeblock}

\begin{codeblock}[title=\textbf{Commandes Feast}]
\begin{minted}{bash}
# Installer Feast
pip install feast

# Initialiser un projet
feast init feature_repo
cd feature_repo

# Appliquer les definitions
feast apply

# Materialiser les features dans le store online
feast materialize 2025-01-01T00:00:00 2025-12-31T23:59:59

# Lister les feature views
feast feature-views list
\end{minted}
\end{codeblock}

\section{Bonnes pratiques de feature engineering}
\label{sec:feature-engineering-bp}

\begin{bonnepratique}[title=\textbf{R\`egles d'or du feature engineering}]
\begin{enumerate}
  \item \textbf{Centraliser les d\'efinitions}~: une seule source de
    v\'erit\'e pour chaque feature (feature store ou module Python d\'edi\'e).
  \item \textbf{Versionner les transformations}~: le code de feature
    engineering est versionn\'e dans Git comme tout autre code.
  \item \textbf{Tester les features}~: v\'erifier les distributions,
    les valeurs manquantes, les bornes attendues.
  \item \textbf{Documenter chaque feature}~: nom, description, unit\'e,
    source, fr\'equence de mise \`a jour.
  \item \textbf{\'Eviter le data leakage}~: ne jamais utiliser d'information
    future lors du calcul des features.
  \item \textbf{Garantir la coh\'erence train/serve}~: utiliser le
    \textbf{m\^eme code} pour calculer les features en entra\^inement et
    en inf\'erence.
\end{enumerate}
\end{bonnepratique}

\begin{antipattern}[title=\textbf{Anti-patterns du feature engineering}]
\begin{itemize}
  \item Calculer des features directement dans le notebook sans les
    centraliser
  \item Utiliser des transformations diff\'erentes en entra\^inement et
    en production (training-serving skew)
  \item Ne pas surveiller la distribution des features en production
    (feature drift)
  \item Cr\'eer des centaines de features sans \'evaluer leur importance
  \item Hardcoder des seuils de discr\'etisation sans les param\'etrer
\end{itemize}
\end{antipattern}

\section{Architecture d'un pipeline de donn\'ees ML}
\label{sec:pipeline-architecture}

Le diagramme suivant montre l'architecture compl\`ete d'un pipeline de
donn\'ees dans un syst\`eme MLOps~:

\begin{center}
\begin{tikzpicture}[
  box/.style={rectangle, draw, rounded corners, minimum width=2.2cm,
    minimum height=0.9cm, align=center, font=\scriptsize},
  arr/.style={-{Stealth[length=2mm]}, thick, steelblue},
  phase/.style={font=\scriptsize\bfseries, text=gray}
]
  % Ingestion
  \node[box, fill=blue!10] (ingest) {Ingestion\\(API, DB, fichiers)};
  \node[phase, above=0.15cm of ingest] {1. Collecte};

  % Validation
  \node[box, fill=blue!15, right=1cm of ingest] (valid) {Validation\\(Great Expect.)};
  \node[phase, above=0.15cm of valid] {2. Validation};

  % Transformation
  \node[box, fill=green!10, right=1cm of valid] (transform) {Transformation\\(nettoyage)};
  \node[phase, above=0.15cm of transform] {3. Transform};

  % Feature eng
  \node[box, fill=green!15, right=1cm of transform] (feat) {Feature\\Engineering};
  \node[phase, above=0.15cm of feat] {4. Features};

  % Feature store
  \node[box, fill=orange!10, right=1cm of feat] (store) {Feature\\Store};
  \node[phase, above=0.15cm of store] {5. Stockage};

  % Training
  \node[box, fill=red!10, below=1.2cm of store] (train) {Entra\^inement};

  % Serving
  \node[box, fill=red!10, below=1.2cm of feat] (serve) {Serving};

  % Arrows
  \draw[arr] (ingest) -- (valid);
  \draw[arr] (valid) -- (transform);
  \draw[arr] (transform) -- (feat);
  \draw[arr] (feat) -- (store);
  \draw[arr] (store) -- (train);
  \draw[arr] (store) -- (serve);

  % Orchestrator
  \node[box, fill=yellow!15, below=1.2cm of transform,
    minimum width=8cm] (orch) {Orchestrateur (Airflow / Prefect / DVC)};
  \draw[dashed, gray] (orch.north) -- (ingest.south);
  \draw[dashed, gray] (orch.north) -- (valid.south);
  \draw[dashed, gray] (orch.north) -- (transform.south);
  \draw[dashed, gray] (orch.north) -- (feat.south);
\end{tikzpicture}
\end{center}

\section{Comparaison des outils d'orchestration}
\label{sec:orchestration-comparison}

\begin{center}
\begin{tabular}{@{}lcccc@{}}
\toprule
\textbf{Crit\`ere} & \textbf{Airflow} & \textbf{Prefect} &
  \textbf{Dagster} & \textbf{Luigi} \\
\midrule
Facilit\'e d'installation & Difficile & Facile & Moyen & Facile \\
Courbe d'apprentissage & Raide & Douce & Moyenne & Douce \\
UI int\'egr\'ee & Oui & Oui & Oui & Basique \\
Gestion des erreurs & Basique & Avanc\'ee & Avanc\'ee & Basique \\
Scalabilit\'e & Tr\`es haute & Haute & Haute & Moyenne \\
Cloud natif & Non & Oui & Non & Non \\
Communaut\'e & Tr\`es large & Croissante & Croissante & Stable \\
Cas d'usage ML & G\'en\'eraliste & ML-friendly & ML-friendly & ETL \\
Typage des donn\'ees & Non & Oui & Oui & Non \\
\bottomrule
\end{tabular}
\end{center}

\begin{attention}[title=\textbf{Choisir le bon outil}]
\begin{itemize}
  \item \textbf{Airflow}~: choix \'eprouv\'e pour les grandes organisations
    avec une \'equipe d'infrastructure d\'edi\'ee.
  \item \textbf{Prefect}~: id\'eal pour les \'equipes ML qui veulent
    d\'emarrer rapidement sans infrastructure lourde.
  \item \textbf{Dagster}~: excellent pour les pipelines de donn\'ees typ\'es
    avec des assets mat\'erialis\'es.
  \item \textbf{Luigi}~: adapt\'e aux pipelines simples et lin\'eaires.
\end{itemize}
\end{attention}

\section{Mini-projet~: pipeline de donn\'ees pour le projet de cours}
\label{sec:mini-project-ch5}

\begin{miniprojet}[title=\textbf{Mini-projet 5~: Pipeline de donn\'ees complet}]
Construisez un pipeline de donn\'ees pour le projet de cours~:
\begin{enumerate}
  \item \textbf{Ingestion}~: \'ecrivez une t\^ache Prefect qui t\'el\'echarge
    un jeu de donn\'ees depuis une URL (par exemple, un dataset UCI ou Kaggle).
  \item \textbf{Validation}~: ajoutez une \'etape qui v\'erifie le sch\'ema
    des donn\'ees (colonnes attendues, types, pas de valeurs aberrantes).
  \item \textbf{Transformation}~: nettoyez les donn\'ees (valeurs manquantes,
    doublons, encodage des variables cat\'egorielles).
  \item \textbf{Feature engineering}~: cr\'eez au moins 5 features
    d\'eriv\'ees pertinentes pour votre probl\`eme.
  \item \textbf{DVC pipeline}~: \'ecrivez un \file{dvc.yaml} qui reproduit
    le m\^eme pipeline avec DVC.
  \item \textbf{Bonus}~: configurez Feast pour stocker et servir vos features.
\end{enumerate}
Livrables~: code source, \file{dvc.yaml}, documentation des features.
\end{miniprojet}

\section{Exercices}
\label{sec:exercises-ch5}

\begin{exercice}[$\bigstar$ --- Pipeline Prefect simple]
\'Ecrivez un pipeline Prefect \`a trois \'etapes (extract, transform, load)
qui~:
\begin{enumerate}
  \item Charge un fichier CSV
  \item Supprime les lignes avec des valeurs manquantes
  \item Sauvegarde le r\'esultat en Parquet
\end{enumerate}
Testez le pipeline localement et v\'erifiez les logs dans le terminal.
\end{exercice}

\begin{exercice}[$\bigstar\bigstar$ --- Pipeline DVC multi-\'etapes]
Cr\'eez un pipeline DVC complet pour un probl\`eme de classification~:
\begin{enumerate}
  \item D\'efinissez 4 \'etapes dans \file{dvc.yaml}~: prepare, featurize,
    train, evaluate
  \item Utilisez \file{params.yaml} pour les hyperparam\`etres
  \item Ex\'ecutez \cli{dvc repro} et v\'erifiez que seules les \'etapes
    modifi\'ees sont r\'e-ex\'ecut\'ees
  \item Comparez les m\'etriques entre deux configurations avec
    \cli{dvc metrics diff}
\end{enumerate}
\end{exercice}

\begin{exercice}[$\bigstar\bigstar\bigstar$ --- Feature store avec Feast]
Mettez en place un feature store complet~:
\begin{enumerate}
  \item Cr\'eez un projet Feast avec au moins deux \texttt{FeatureView}
  \item D\'efinissez des features pour deux entit\'es diff\'erentes (par
    exemple, utilisateur et produit)
  \item Mat\'erialisez les features dans le store online
  \item \'Ecrivez un script d'entra\^inement qui r\'ecup\`ere les features
    historiques via \cli{get\_historical\_features}
  \item \'Ecrivez un script de serving qui r\'ecup\`ere les features en
    temps r\'eel via \cli{get\_online\_features}
  \item V\'erifiez que les features sont identiques entre entra\^inement
    et serving (pas de training-serving skew)
\end{enumerate}
\end{exercice}

\section{Aide-m\'emoire}

\begin{cheatsheet}[title=\textbf{R\'esum\'e du chapitre 5}]
\begin{tabular}{@{}p{3.5cm}p{9cm}@{}}
\toprule
\textbf{Concept} & \textbf{Point cl\'e} \\
\midrule
Pipeline de donn\'ees & S\'equence automatis\'ee~: ingestion $\to$
  validation $\to$ transformation $\to$ features \\
ETL vs ELT & ETL~: transformer avant le chargement~; ELT~: transformer
  apr\`es \\
Prefect & \cli{@flow}, \cli{@task}, retries, UI int\'egr\'ee \\
DVC pipelines & \file{dvc.yaml}, \cli{dvc repro}, ne r\'e-ex\'ecute que
  le n\'ecessaire \\
Feature store & Centralise les features, garantit la coh\'erence
  train/serve \\
Feast & \cli{feast apply}, \cli{get\_historical\_features},
  \cli{get\_online\_features} \\
Training-serving skew & Diff\'erence entre features en entra\^inement et
  en production \\
\bottomrule
\end{tabular}
\end{cheatsheet}
